\section{Versiones y evidencia de entrega}\label{sec:versiones-entrega}

La entrega reúne la base integrada de CPU, UART, Debug, plataforma y software
de PC, y el aporte de aceptación de la \href{\projecturl/pull/37}{PR N.º~37}.
La tabla~\ref{tab:revisiones-entrega} fija las revisiones de referencia. Los
hashes abreviados enlazan a sus commits completos.

\begin{table}[H]
  \centering
  \small
  \begin{tabularx}{\textwidth}{L{4.0cm}L{3.1cm}Y}
    \toprule
    \textbf{Componente} & \textbf{Revisión} & \textbf{Identificación} \\
    \midrule
    Base integrada &
    \href{\projecturl/commit/\baserevision}{\code{7025925f949c}} &
    RTL, contratos y \tech{viewrv} 0.7.0. \\ \addlinespace
    Build de FPGA &
    \href{\projecturl/commit/d2ef373194b396f5a4367af22fe7e243ada5d49e}{\code{d2ef373194b3}} &
    Vivado 2025.2, perfil \code{REFERENCE}; fuentes equivalentes a la base. \\
    \addlinespace
    Aceptación de la PR N.º~37 &
    \href{\projecturl/commit/\acceptancerevision}{\code{cb3270b19b27}} &
    Ejecutor, pruebas y guía de los 151 IDs. \\
    \bottomrule
  \end{tabularx}
  \caption{Revisiones de referencia de la entrega.}
  \label{tab:revisiones-entrega}
\end{table}

La \href{\projecturl/releases/tag/tp3-17-evidencia}{publicación de evidencia}
ofrece dos archivos:

\begin{itemize}
  \item \evidencefile{resultados-por-id.csv}{\code{resultados-por-id.csv}}:
    151 filas, con 150 aprobaciones en placa y una excepción física,
    \code{TEST-UART-013/C009}, aprobada en simulación de sistema.
  \item \evidencefile{tp3-issue-17-evidencia.zip}{\code{tp3-issue-17-evidencia.zip}}:
    capturas por ID, logs de \tech{viewrv}, programas, resultados y el
    directorio \code{vivado/}, con bitstream, manifiesto y reportes.
\end{itemize}

Dentro del ZIP, \code{uart-regression-7025925.json} contiene las 44 ejecuciones
de sistema aprobadas, incluida C009. \code{viewrv-comparison.json} resume los
resultados del recorrido CLI. Los directorios de los intentos de botón
conservan los intentos iniciales y sus repeticiones; el CSV consolida el estado
final. La huella del bitstream figura en el
capítulo~\ref{chap:plataforma-timing}.


\section{Entorno de RTL y assembler}

Las comprobaciones del repositorio se ejecutan en la imagen publicada
\code{tp3-dev} de GitHub Container Registry, fijada por digest en
\code{tp3/docker/image.ref}. La imagen contiene Icarus Verilog 13.0, Verilator
5.020, Verible, check-jsonschema, Ruff, ShellCheck, shfmt, GNU Make y CPython
3.12.3. La descripción completa y las entradas fijadas se encuentran en
\doc{development_environment.md}{Entorno de compilación y simulación}.

Desde la raíz del repositorio, en una terminal Bash:

\begin{codebox}[language=bash]
TP3_IMAGE="$(cat tp3/docker/image.ref)"
docker pull "$TP3_IMAGE"
docker run --rm --network none \
  --user "$(id -u):$(id -g)" \
  --mount "type=bind,src=$PWD,dst=/workspace" \
  "$TP3_IMAGE" make check
\end{codebox}

La descarga requiere acceso al registro. Las comprobaciones se ejecutan sin
red. La imagen aporta las herramientas; el RTL, los bancos y las pruebas se
toman del checkout montado.

\section{Comprobaciones desde \code{tp3}}

Dentro del contenedor, desde \code{/workspace/tp3}:

\begin{codebox}[language=bash]
make help
make check
\end{codebox}

\code{make check} ejecuta:

\begin{enumerate}
  \item \code{format-check}: formato SystemVerilog, Python y Bash;
  \item \code{lint}: Verible, Verilator, Ruff y ShellCheck;
  \item \code{validate}: esquemas JSON y catálogo de verificación;
  \item \code{review-cpu}: estructura del CPU en tres configuraciones;
  \item \code{test-infra}: pruebas de la infraestructura;
  \item \code{test-software}: las 114 pruebas del assembler;
  \item \code{test-rtl}: simulaciones RTL del manifiesto.
\end{enumerate}

Los comandos parciales son:

\begin{codebox}[language=bash]
make test-infra
make test-software
make test-rtl
make test-rtl TESTS="TEST-PIPE-005/C001"
\end{codebox}

La guía de selección, los filtros y el formato JSON se consultan en
\doc{verification/test_selection.md}{Ejecución y selección de pruebas}.

Los targets guardan sus resultados en \code{build/reports/}. Para conservar
también la salida completa de la regresión, desde \code{tp3}, en Bash:

\begin{codebox}[language=bash]
mkdir -p build/reports
set -o pipefail
make check 2>&1 | tee build/reports/check.log
\end{codebox}

\section{Entorno y pruebas de viewrv}

El entorno de la aplicación es independiente del contenedor RTL. La versión
0.7.0 fija CPython 3.14.9 y Node.js 24.22.1, con dependencias en
\code{requirements-linux.lock} y \code{package-lock.json}. En Linux, desde
\code{tp3}, con esas versiones disponibles:

\begin{codebox}[language=bash]
python3.14 -m venv software/.venv
software/.venv/bin/python -m pip install --require-hashes \
  -r software/requirements-linux.lock
cd software/gui
npm ci --no-audit --no-fund
node node_modules/electron/install.js
npm run check
npm run test:core
npm run test:gui
\end{codebox}

La instalación descarga dependencias. Una vez preparadas, las pruebas utilizan
recursos locales. \code{check} comprueba formato, compila TypeScript y ejecuta
las pruebas unitarias de GUI; \code{test:core} ejecuta la suite Python de core,
CLI y protocolo; \code{test:gui} recorre las ventanas sobre la construcción
generada. Este último requiere una sesión gráfica X11 en Linux y conserva
resultados bajo \code{software/build/gui-check/}. \code{npm run check:all}
reúne los tres niveles.

Los adaptadores controlados y pseudo-terminales permiten comprobar el cliente
sin FPGA. Esas pruebas no acreditan por sí solas la GUI sobre Basys 3 ni la
distribución Windows. Los procedimientos de aplicación se detallan en el
\href{\projecturl/blob/\baserevision/tp3/software/README.md}{README de viewrv}.


\section{Resultados de la campaña registrada}

La campaña documentada en el
\href{\projecturl/blob/\baserevision/tp3/reports/issue-16.md}{informe de integración y timing}
se ejecutó sin red. La corrida principal usó Icarus
Verilog 13.0 y Verilator 5.054; la repetición de los tres casos de recuperación
usó Verilator 5.020, la versión fijada por el repositorio. Los resultados
fueron:

\begin{table}[H]
  \centering
  \small
  \begin{tabular}{lr}
    \toprule
    \textbf{Comprobación} & \textbf{Resultado} \\
    \midrule
    Formato, lint, validación y revisión del CPU & OK (sin errores) \\
    Infraestructura & OK \\
    Assembler & 114 pruebas, OK \\
    Simulaciones RTL & 2695 ejecuciones: 2692 PASS, 0 FAIL, 3 ERROR \\
    Repetición con Verilator 5.020 & 3 de 3 PASS \\
    Casos con estado pendiente en el runner & 4 \\
    \bottomrule
  \end{tabular}
  \caption{Resultado de la regresión ejecutada.}
  \label{tab:corrida}
\end{table}

Los tres errores correspondieron a las ejecuciones de recuperación
\code{TEST-UART-006/C001}, \code{C013} y \code{C014} con Verilator 5.054. Las
tres pasaron al repetirse con Verilator 5.020, la versión fijada por el
repositorio. La diferencia se conserva como una dependencia de la versión del
simulador y no como un fallo funcional del diseño. El resto de la regresión no
presentó fallos ni diferencias de resultado.

Los registros de ambas corridas se identifican como \code{regression.json} y
\code{retry-5.020.json} en el informe de integración. Una nueva ejecución de los
comandos anteriores genera su propia evidencia bajo \code{build/reports/}; no
se la confunde con los resultados históricos de la entrega.

\section{Selección de casos}

Para ejecutar un grupo o un caso:

\begin{codebox}[language=bash]
bash scripts/run_tests.sh --test TEST-PIPE-005 --list
bash scripts/run_tests.sh --test TEST-PIPE-005/C001 --format json
bash scripts/run_tests.sh --test TEST-EXEC-001 --test TEST-MEM-003
bash scripts/run_tests.sh --all --format json
\end{codebox}

\code{--test} es repetible y se puede combinar con \code{--list} para
inspeccionar sin ejecutar. \code{--all} ejecuta todos los casos implementados e
informa los pendientes. Pedir individualmente un caso pendiente es un error;
seleccionar un grupo informa sus pendientes sin ejecutarlos.

\section{Aceptación en placa}

Los comandos de esta sección requieren el aporte de la PR N.º~37, incorporado
en el checkout o disponible en la revisión \code{cb3270b19b27}. No están
incluidos en la base \code{7025925f949c} por sí sola. Para obtener una copia
independiente y fijada, desde una carpeta de trabajo:

\begin{codebox}[language=bash]
git clone https://github.com/FrancoPerotti/computer-architecture-tps.git \
  tp3-aceptacion
cd tp3-aceptacion
git checkout --detach cb3270b19b2721f3e72a23ab23e333b394d113b9
cd tp3
\end{codebox}

Con Python 3.11 o posterior, preparar los planes, comprobar el transporte y
ejecutar las pruebas del host:

\begin{codebox}[language=bash]
python3 -m tools.board_acceptance --output build/board-acceptance/prepared
python3 -m tools.board_acceptance --simulate \
  --output build/board-acceptance/transport
python3 -m unittest -v scripts.tests.test_board_acceptance
\end{codebox}

La preparación exporta 151 planes: 150 preparados y C009 pendiente de
instrumentación. \code{--simulate} ejecuta únicamente \code{smoke} con un
transcript fijo, sin abrir UART ni simular la CPU. Los nueve tests del host
comprueban selección, framing, comparación de snapshots y recuperación.
Las proyecciones sobre \code{processor_top} se ejecutan por separado:

\begin{codebox}[language=bash]
python3 -m tools.board_acceptance.check_cpu \
  --output build/board-acceptance/rtl
\end{codebox}

Ese comando requiere Icarus y las herramientas del entorno RTL. Para ensayar
con la placa se necesita, además, pySerial 3.5, el bitstream \code{REFERENCE}
ya programado y el puerto libre. En la misma terminal Bash, desde \code{tp3},
asignar las rutas reales del puerto y del archivo descargado:

\begin{codebox}[language=bash]
TP3_UART="/dev/ttyUSB1"
TP3_BITSTREAM="/ruta/a/evidencia/vivado/tp3-reference.bit"
TP3_RTL_SHA="d2ef373194b396f5a4367af22fe7e243ada5d49e"
python3 -m tools.board_acceptance --case smoke \
  --port "$TP3_UART" --bitstream "$TP3_BITSTREAM" \
  --rtl-commit "$TP3_RTL_SHA" --operators Franco Federico \
  --notes "Basys 3; identificar placa y turno de operadores" \
  --output build/board-acceptance/plate-smoke
\end{codebox}

El ejecutor registra el bitstream, pero no programa la FPGA. La programación
previa se realiza con \tech{viewrv}; cerrar su conexión UART antes del ensayo
directo. Para el lote automático, omitir \code{--case smoke} y utilizar otro
directorio de salida. Los casos de botón requieren \code{--manual} y el ID
exacto con \code{--case}; C009 conserva el procedimiento de instrumentación
específico.

El ejecutor abre el puerto con exclusión en Linux y usa \num{19200} baud, 8N1 y
sin control de flujo. Cada ejecución exige un directorio nuevo y conserva
\code{session.json}, \code{status.json} y los archivos \code{plan.json},
\code{uart.jsonl} y \code{result.json} por caso. La
\acceptancedoc{guía de aceptación de la PR N.º~37} documenta los turnos, los
estímulos de BTNC y el recorrido CLI de \tech{viewrv}. Guardar el bitstream y
los reportes de Vivado junto con el ensayo.
